%%%%%%%%%%%%%%%%%%%%%%%%%%%%%%%%%%%%%%%%%
% Developer CV
% LaTeX Template
% Version 1.0 (28/1/19)
%
% This template originates from:
% http://www.LaTeXTemplates.com
%
% Authors:
% Jan Vorisek (jan@vorisek.me)
% Based on a template by Jan Küster (info@jankuester.com)
% Modified for LaTeX Templates by Vel (vel@LaTeXTemplates.com)
%
% License:
% The MIT License (see included LICENSE file)
%
%%%%%%%%%%%%%%%%%%%%%%%%%%%%%%%%%%%%%%%%%

%----------------------------------------------------------------------------------------
%	PACKAGES AND OTHER DOCUMENT CONFIGURATIONS
%----------------------------------------------------------------------------------------

\documentclass[9pt]{developercv} % Default font size, values from 8-12pt are recommended

\usepackage{booktabs}
\usepackage{makecell}
\setlength\LTleft{0pt}
\setlength\LTright{0pt}
\def\tightlist{}
\providecommand{\tightlist}{%
  \setlength{\itemsep}{0pt}\setlength{\parskip}{0pt}}
%\usepackage{titlesec}
%
%\titleformat{\paragraph}[leftmargin]
%{\normalfont
%\sffamily\bfseries\filleft}
%{}{0pt}{}
%\titlespacing{\paragraph}
%{4pc}{1.5ex plus .1ex minus .2ex}{0.5pc}


%----------------------------------------------------------------------------------------

\begin{document}

%----------------------------------------------------------------------------------------
%	TITLE AND CONTACT INFORMATION
%----------------------------------------------------------------------------------------

\begin{minipage}[t]{0.45\textwidth} % 45% of the page width for name
	\vspace{-\baselineskip} % Required for vertically aligning minipages
	
	% If your name is very short, use just one of the lines below
	% If your name is very long, reduce the font size or make the minipage wider and reduce the others proportionately
	\colorbox{black}{{\HUGE\textcolor{white}{\textbf{\MakeUppercase{Pawan}}}}} % First name
	
	\colorbox{black}{{\HUGE\textcolor{white}{\textbf{\MakeUppercase{Bharadwaj}}}}} % Last name
	
	\vspace{6pt}
	
	{\huge Research Scientist (Applied Mathematics \& Geophysics)} % Career or current job title
\end{minipage}
\begin{minipage}[t]{0.275\textwidth} % 27.5% of the page width for the first row of icons
	\vspace{-\baselineskip} % Required for vertically aligning minipages
	
	% The first parameter is the FontAwesome icon name, the second is the box size and the third is the text
	% Other icons can be found by referring to fontawesome.pdf (supplied with the template) and using the word after \fa in the command for the icon you want
	\icon{MapMarker}{12}{Cambridge, MA, USA}\\
	\icon{Phone}{12}{+1 8573205756}\\
	\icon{At}{12}{\href{mailto:pawbz@mit.edu}{pawbz@mit.edu}}\\	
\end{minipage}
\begin{minipage}[t]{0.275\textwidth} % 27.5% of the page width for the second row of icons
	\vspace{-\baselineskip} % Required for vertically aligning minipages
	
	% The first parameter is the FontAwesome icon name, the second is the box size and the third is the text
	% Other icons can be found by referring to fontawesome.pdf (supplied with the template) and using the word after \fa in the command for the icon you want
	\icon{Globe}{12}{\href{https://pawbz.github.io}{pawbz.github.io}}\\
	\icon{Github}{12}{\href{https://github.com/pawbz}{github.com/pawbz}}\\
	\icon{Google}{12}{\href{https://scholar.google.co.in/citations?user=EMbC7l8AAAAJ&hl=en}{Google Scholar}}\\
\end{minipage}

\vspace{0.5cm}

%----------------------------------------------------------------------------------------
%	INTRODUCTION, SKILLS AND TECHNOLOGIES
%----------------------------------------------------------------------------------------

\cvsect{Who Am I?}

\begin{minipage}[t]{0.4\textwidth} % 40% of the page width for the introduction text
	\vspace{\baselineskip} % Required for vertically aligning minipages
	A driven and knowledgeable researcher offering proven background in applied mathematics and geophysics. History of developing novel algorithms related geophysical inverse problems, signal processing and machine learning.
	Authored more than 10 high-impact publications and patents.
%A researcher 
%with an expertise in inverse problems and 
%signal processing relevant to wave propagation and 
%scattering --- passionate to solve real-life challenges pertaining to seismic, radar and medical imaging. 

%Recently working on projects that bridge 
%	the ideas of machine learning to the physics of wave propagation.
\end{minipage}
\hfill % Whitespace between
\begin{minipage}[t]{0.5\textwidth} % 50% of the page for the skills bar chart
	\vspace{-\baselineskip} % Required for vertically aligning minipages
	\begin{barchart}{5.5}
		\baritem{Seismology}{47}
		\baritem{Signal Processing}{50}
		\baritem{Machine Learning}{40}
		%\baritem{PDE-constrained Optimization}{40}
		\baritem{Optimization and Inverse Problems}{45}
		\baritem{High-performance Computing}{38}
		%\baritem{Seismic Interferometry}{45}
		\baritem{Interactive Teaching/ Learning}{38}
		\baritem{Public Speaking}{40}
	\end{barchart}
\end{minipage}

\begin{center}
	\bubbles{5/Julia, 4/Jupyter, 3/bash, 4/git, 3/Python, 4/FORTRAN}
\end{center}

\cvsect{Research
Positions}\label{research-positions}

\paragraph{12/2016---11/2019}\label{section}

\emph{Postdoctoral
Associate},
Department
of
Mathematics
\& Earth
Resources
Laboratory,
Massachusetts
Institute
of
Technology
(MIT),
Cambridge,
USA.

\begin{itemize}
\item
  Furthered
  field
  knowledge
  in
  machine
  learning
  by
  developing
  data
  processing
  methods
  such as
  blind
  deconvolution
  and
  source
  separation.
\item
  Research
  demonstrated
  the
  utility
  of these
  methods
  in
  geophysics
  e.g.,
  seismic
  waves
  recorded
  at
  \textasciitilde{}500
  stations
  worldwide
  were
  factored
  to
  automatically
  extract
  key
  features
  pertaining
  to
  earthquake
  characterization.
\item
  Presented
  research
  results
  at
  applied
  mathematics,
  seismology
  and
  applied
  geophysics
  conferences
  to
  audience
  of 100+,
  and
  published
  the
  results
  in major
  IEEE and
  seismology
  journals.
\item
  Developed
  HPC-optimized
  Julia
  open-source
  software
  libraries
  to
  perform
  regression
  during
  the
  above-mentioned
  processing.
\item
  Gathered
  and
  analyzed
  data on
  research
  interests
  of
  colleagues
  in order
  to
  successfully
  organize
  departmental
  seminars
  for two
  semesters.
\end{itemize}

\paragraph{09/2012---11/2016}\label{section-1}

\emph{Graduate
Researcher},
Civil
Engineering
and
Geosciences,
Delft
University
of
Technology
(TU
Delft),
Delft, The
Netherlands.

\begin{itemize}
\item
  Developed
  computational
  imaging
  algorithms
  for a
  seismic
  system
  that
  performs
  real-time
  (within
  \textasciitilde{}5-7
  minutes)
  mapping
  of the
  Earth's
  structure
  ahead of
  a tunnel
  boring
  machine.~
\item
  Successfully
  coordinated
  with a
  group of
  civil
  and
  mechanical
  engineers
  to test
  the
  utility
  of this
  system
  for
  hazard
  assessment
  during
  excavation.
\item
  Presented
  at
  industry
  conferences
  to
  audiences
  of up to
  50-100
  professionals.
\item
  Engaged
  \textasciitilde{}70-80
  students
  in
  several
  workshops
  on
  seismic
  inverse
  problems/
  data
  analysis
  that
  focus on
  teaching
  via
  interactive
  computing
  with
  high-level
  programming
  languages.
\item
  Effectively
  maintained
  and
  tested
  departmental
  software
  libraries.
\end{itemize}

\paragraph{05/2010---12/2012}\label{section-2}

Multiple
short-term
visits as
a
\emph{Research
Assistant}
in
Physical
Science
and
Engineering,
King
Abdullah
University
of Science
and
Technology
(KAUST),
Saudi
Arabia.

\begin{itemize}
\item
  Spearheaded
  the
  development
  super-virtual
  interferometry
  (SVI),
  which
  performs
  denoising
  and
  thereby
  extracts
  coherent
  energy
  from
  seismic
  signals.
\item
  SVI is
  now
  routinely
  used
  both in
  academia
  and the
  energy
  industry
  to
  process
  measurements
  for
  enhanced
  subsurface
  seismic
  imaging.
\end{itemize}

\cvsect{Education}\label{education}

\paragraph{2012---2016}\label{section-3}

\emph{PhD
in
Geophysics},
Delft
University
of
Technology
(TU
Delft),
Delft, The
Netherlands,
Thesis:
Seismic
Waveform
Inversion:
Bump
Functional,
Parameterization
Analysis
and
Imaging
Ahead of a
Tunnel-Boring
Machine.
\href{https://repository.tudelft.nl/islandora/object/uuid:8d59fd51-cbac-408d-a3e0-6a15891c2695/datastream/OBJ/download}{{[}pdf{]}}

\paragraph{2007---2012}\label{section-4}

\emph{Master
of Science
and
Technology,
Applied
Geophysics},
five year
integrated
degree in
Indian
Institute
of
Technology
(Indian
School of
Mines),
IIT (ISM),
Dhanbad,
India.

\cvsect{Teaching
Experience
and
Certification}\label{teaching-experience-and-certification}

\paragraph{Spring
2019}\label{spring-2019}

\emph{Kaufman
Teaching
Certificate
Program
(KTCP)}.
Demonstrated
a strong
commitment
to the
teaching
enterprise
by
completing
several
workshops
and
contributing
to class
discussions
and
activities
designed
to help
practice
teaching
methods in
Massachusetts
Institute
of
Technology.

\paragraph{Fall
2013/2014}\label{fall-20132014}

\emph{Computational
Geophysics
Lab
Instructor}.
Engaged
students
in several
workshops
that focus
on
teaching
via
interactive
computing,
Delft
University
of
Technology.

\cvsect{Awards
and
Grants}\label{awards-and-grants}

\paragraph{2018}\label{section-5}

\emph{Best
Paper
Award}
during the
annual TU
Delft
Geosciences
Research
Meeting.

\paragraph{2014---2016}\label{section-6}

\emph{Computational
Resource
Grant} to
develop
near-surface
seismic
imaging
algorithms
for
real-time
risk
assessment
during
tunnel
excavation,
Netherlands
Organization
for
Scientific
Research
(NWO).

\paragraph{2008---2012}\label{section-7}

Innovation
in Science
Pursuit
for
Inspired
Research
(INSPIRE)
Scholarship,
Department
of Science
and
Technology,
India.

\paragraph{2008}\label{section-8}

Scholarship
for Higher
Education,
Indian
School of
Mines.

\cvsect{Invited
Talks}\label{invited-talks}

\paragraph{11/2019}\label{section-9}

King
Abdullah
University
of Science
and
Technology
(KAUST),
Saudi
Arabia.

\paragraph{06/2019}\label{section-10}

Woods Hole
Oceanographic
Institution,
MA, USA.

\paragraph{04/2019}\label{section-11}

Schlumberger-Doll
Research
Center,
Cambridge,
USA.

\paragraph{01/2019}\label{section-12}

Center for
Earth
Sciences,
Indian
Institute
of Science
(IISc).

\paragraph{10/2018}\label{section-13}

Caltech
Seismological
Laboratory,
California
Institute
of
Technology.

\paragraph{07/2016}\label{section-14}

Earth
Resources
Laboratory,
Massachusetts
Institute
of
Technology.

\paragraph{05/2016}\label{section-15}

Department
of Earth
Sciences
in Utrecht
University,
NL.

\paragraph{04/2016}\label{section-16}

Institute
of
Geophysics
and
Planetary
Physics in
University
of
California,
Santa
Cruz, USA.

\cvsect{Short-term
Academic
Visits}\label{short-term-academic-visits}

\paragraph{10/2019}\label{section-17}

Visiting
Research
Scientist,
Department
of Marine
Geosciences,
Institut
de
Physique
du Globe
de Paris.

\paragraph{09/2011}\label{section-18}

Guest
Researcher,
Institute
for
Geophysics,
University
of Texas
at Austin,
USA.

\cvsect{Service}\label{service}

\begin{itemize}
\tightlist
\item
  Reviewer:

  \begin{itemize}
  \tightlist
  \item
    Geophysical
    Journal
    International,
    Geophysics,
    IEEE
    Transactions
    on
    Geoscience
    and
    Remote
    Sensing.
  \end{itemize}
\item
  Professional
  Associations:

  \begin{itemize}
  \tightlist
  \item
    American
    Geophysical
    Union
    (AGU),
    Society
    of
    Exploration
    Geophysicists
    (SEG),
    European
    Association
    of
    Geoscientists
    and
    Engineers
    (EAGE),
    Society
    for
    Industrial
    and
    Applied
    Mathematics
    (SIAM).
  \end{itemize}
\end{itemize}

\cvsect{Peer-reviewed
Journal
Articles}\label{peer-reviewed-journal-articles}

\paragraph{2019}\label{section-19}

{[}9{]}
\textbf{\emph{Pawan
Bharadwaj}},
Chunfang
Meng,
Michael
Fehler,
Laurent
Demanet
and Aimé
Fournier,
Redshift
of
earthquakes
via
focused
blind
deconvolution
of
teleseisms
(in
review).

\paragraph{2018}\label{section-20}

{[}8{]}
\textbf{\emph{Pawan
Bharadwaj}},
Laurent
Demanet,
Aimé
Fournier,
Focused
blind
deconvolution,
\texttt{IEEE\ Transactions\ on\ Signal\ Processing}.
\href{https://github.com/pawbz/pawbz.github.io/blob/src/pdf/papers/Bharadwaj_FBD_arXiv1.pdf}{{[}pdf{]}}

\paragraph{2017}\label{section-21}

{[}7{]}
\textbf{\emph{Pawan
Bharadwaj}},
Wim
Mulder,
Guy
Drijkoningen,
Jan
Thorbecke,
Boris
Neducza,
and Rob
Jenneskens,
A
shear-wave
seismic
system
using full
waveform
inversion
to look
ahead of a
tunnel-boring
machine,
\texttt{Near\ Surface\ Geophysics}.
\href{https://github.com/pawbz/pawbz.github.io/blob/src/pdf/papers/NSG15_Bharadwaj.pdf}{{[}pdf{]}}

\paragraph{2016}\label{section-22}

{[}6{]}
\textbf{\emph{Pawan
Bharadwaj}},
Wim
Mulder,
and Guy
Drijkoningen,
Full
waveform
inversion
with an
auxiliary
bump
functional,
\texttt{Geophysical\ Journal\ International}.
\href{https://github.com/pawbz/pawbz.github.io/blob/src/pdf/papers/FWI_bump.pdf}{{[}pdf{]}}

\paragraph{2014}\label{section-23}

{[}5{]}
\textbf{\emph{Pawan
Bharadwaj}},
Tarje
Nissen-Meyer,
Gerard
Schuster,
and Martin
Mai,
Enhancing
core-diffracted
arrivals
by
supervirtual
interferometry,
\texttt{Geophysical\ Journal\ International}.
\href{https://github.com/pawbz/pawbz.github.io/blob/src/pdf/papers/GJI_CMB.pdf}{{[}pdf{]}}

\paragraph{2013}\label{section-24}

{[}4{]}
\textbf{\emph{Pawan
Bharadwaj}},
Xin Wang,
Gerard
Schuster,
and Kirk
McIntosh,
Increasing
the number
and
signal-to-noise
ratio of
OBS traces
with
supervirtual
refraction
interferometry
and
free-surface
multiples,
\texttt{Geophysical\ Journal\ International}.
\href{https://github.com/pawbz/pawbz.github.io/blob/src/pdf/papers/svi_GJI2.pdf}{{[}pdf{]}}

\paragraph{2012}\label{section-25}

{[}3{]}
\textbf{\emph{Pawan
Bharadwaj}},
Gerard
Schuster,
Ian
Mallinson,
and Wei
Dai,
Theory of
supervirtual
refraction
interferometry,
\texttt{Geophysical\ Journal\ International}.
\href{https://github.com/pawbz/pawbz.github.io/blob/src/pdf/papers/svi_GJI1.pdf}{{[}pdf{]}}

\paragraph{}\label{section-26}

{[}2{]}
Vinay
Kumar
Srivastava,
Devendra
Nath Giri,
and
\textbf{\emph{Pawan
Bharadwaj}},
Study and
mapping of
ground
water
prospect
using
remote
sensing,
GIS and
geoelectrical
resistivity
techniques--a
case study
of Dhanbad
district,
Jharkhand,
India,
\texttt{Journal\ of\ Indian\ Geophysical\ Union}.

\paragraph{2011}\label{section-27}

{[}1{]}
Ian
Mallinson,
\textbf{\emph{Pawan
Bharadwaj}},
Gerard
Schuster,
and Helmut
Jakubowicz,
Enhanced
refractor
imaging by
supervirtual
interferometry,
\texttt{The\ Leading\ Edge}.
\href{https://github.com/pawbz/pawbz.github.io/blob/src/pdf/papers/svi_TLE.pdf}{{[}pdf{]}}

\cvsect{Peer-reviewed
Conference
Articles
\&
Presentations}\label{peer-reviewed-conference-articles-presentations}

\paragraph{2019}\label{section-28}

{[}9{]}
Aimé
Fournier,
Charles­Henri
Clerget,
\textbf{\emph{Pawan
Bharadwaj}},
A
seismoelectric
inverse
problem
with
well-log
data and
borehole-confined
acquisition,
Society of
Exploration
Geophysicists,
\texttt{SEG\ Annual\ Meeting}.

\paragraph{2018}\label{section-29}

{[}8{]}
\textbf{\emph{Pawan
Bharadwaj}},
Laurent
Demanet,
Aimé
Fournier,
Focused
blind
deconvolution
of
interferometric
Green's
functions,
Society of
Exploration
Geophysicists,
\texttt{SEG\ Annual\ Meeting}.
\href{https://github.com/pawbz/pawbz.github.io/blob/src/pdf/abstracts/Pawan_FIBD_SEG.pdf}{{[}pdf{]}}

\paragraph{2017}\label{section-30}

{[}7{]}
\textbf{\emph{Pawan
Bharadwaj}},
Laurent
Demanet,
Aimé
Fournier,
Deblending
random
seismic
sources
via
independent
component
analysis,
Society of
Exploration
Geophysicists,
\texttt{SEG\ Annual\ Meeting}.
\href{https://github.com/pawbz/pawbz.github.io/blob/src/pdf/abstracts/Pawan_ICA_SEG.pdf}{{[}pdf{]}}

\paragraph{2016}\label{section-31}

{[}6{]}
\textbf{\emph{Pawan
Bharadwaj}},
Guy
Drijkoningen,
Wim
Mulder,
Thomas
Tscharner,
and Rob
Jenneskens,
A
shear-wave
seismic
system to
look ahead
of a
tunnel
boring
machine,
Proceedings
of the
Society
for
Mining,
Metallurgy
\&
Exploration,
\texttt{World\ Tunneling\ Congress}.
\href{https://github.com/pawbz/pawbz.github.io/blob/src/pdf/abstracts/WTC2016_Pawan.pdf}{{[}pdf{]}}

\paragraph{2015}\label{section-32}

{[}5{]}
\textbf{\emph{Pawan
Bharadwaj}},
Wim
Mulder,
Guy
Drijkoningen,
and Ralph
Reijnen,
Looking
ahead of a
tunnel
boring
machine
with 2-D
SH full
waveform
inversion,
European
Association
of
Geoscientists
\&
Engineers
Conference
and
Exhibition,
\texttt{EAGE\ Annual\ Meeting}.
\href{https://github.com/pawbz/pawbz.github.io/blob/src/pdf/abstracts/Pawan_looking_aheadEAGE.pdf}{{[}pdf{]}}

\paragraph{}\label{section-33}

{[}4{]}
\textbf{\emph{Pawan
Bharadwaj}},
Wim
Mulder,
and Guy
Drijkoningen,
An
auxiliary
bump
functional
to
overcome
cycle
skipping
in
waveform
inversion,
European
Association
of
Geoscientists
\&
Engineers
Conference
and
Exhibition-Workshops,
\texttt{EAGE\ Annual\ Meeting}.
\href{https://github.com/pawbz/pawbz.github.io/blob/src/pdf/abstracts/bump_functional_workshop.pdf}{{[}pdf{]}}

\paragraph{2014}\label{section-34}

{[}3{]}
\textbf{\emph{Pawan
Bharadwaj}},
Wim
Mulder,
and Guy
Drijkoningen,
Parameterization
for 2-D SH
full
waveform
inversion,
European
Association
of
Geoscientists
\&
Engineers
Conference
and
Exhibition,
\texttt{EAGE\ Annual\ Meeting}.
\href{https://github.com/pawbz/pawbz.github.io/blob/src/pdf/abstracts/parameterization_EAGE.pdf}{{[}pdf{]}}

\paragraph{2013}\label{section-35}

{[}2{]}
\textbf{\emph{Pawan
Bharadwaj}},
Wim
Mulder,
and Guy
Drijkoningen,
Multi-objective
full
waveform
inversion
in the
absence of
low
frequencies,
Society of
Exploration
Geophysicists,
\texttt{SEG\ Annual\ Meeting}.
\href{https://github.com/pawbz/pawbz.github.io/blob/src/pdf/abstracts/Multi_objectiveFWI.pdf}{{[}pdf{]}}

\paragraph{2011}\label{section-36}

{[}1{]}
\textbf{\emph{Pawan
Bharadwaj}},
Gerard
Schuster,
and Ian
Mallinson,
Super-virtual
refraction
interferometry:
theory,
Society of
Exploration
Geophysicists,
\texttt{SEG\ Annual\ Meeting}.
\href{https://github.com/pawbz/pawbz.github.io/blob/src/pdf/abstracts/SEG_SVITheory.pdf}{{[}pdf{]}}

\cvsect{Conference
Abstracts}\label{conference-abstracts}

\paragraph{2018}\label{section-37}

{[}5{]}
\textbf{\emph{Pawan
Bharadwaj}},
Yusuke
Mukuhira,
Laurent
Demanet,
Michael
Fehler,
Aimé
Fournier,
Focused
blind
deconvolution
of
earthquake
signals,
American
Geophysical
Union,
\texttt{AGU\ Fall\ Meeting}.

\paragraph{}\label{section-38}

{[}4{]}
\textbf{\emph{Pawan
Bharadwaj}},
Laurent
Demanet,
Aimé
Fournier,
Sparsity
in
multichannel
blind
deconvolution
via
focusing
constraints,
\texttt{SIAM\ Conference\ on\ Mathematical\ \&\ Computational\ Issues\ in\ the\ Geosciences}.

\paragraph{2012}\label{section-39}

{[}3{]}
Gerard
Schuster,
\textbf{\emph{Pawan
Bharadwaj}},
Tarje
Nissen-Meyer,
and Martin
Mai,
Extracting
CMB
diffractions
by
super-virtual
interferometry,
American
Geophysical
Union,
\texttt{AGU\ Fall\ Meeting}.

\paragraph{2011}\label{section-40}

{[}2{]}
Gerard
Schuster,
\textbf{\emph{Pawan
Bharadwaj}},
and Kirk
McIntosh,
Refraction
interferometry
and
free-surface
multiples
for
large-offset
OBS data,
American
Geophysical
Union,
\texttt{AGU\ Fall\ Meeting}.

\paragraph{2010}\label{section-41}

{[}1{]}
Gerard
Schuster
and
\textbf{\emph{Pawan
Bharadwaj}},
Extending
the
aperture
and
enhancing
the S/N
ratio of
refraction
imaging by
super-virtual
interferometry,
American
Geophysical
Union,
\texttt{AGU\ Fall\ Meeting}.

\cvsect{Patents}\label{patents}

\paragraph{2019}\label{section-42}

Aimé
Fournier,
\textbf{\emph{Pawan
Bharadwaj}},
Laurent
Demanet,
Sequential
estimation
while
drilling,
International
Patent
Application
PCT/US2019/054945,
U.S.
Patent
Application
62/857971.

\paragraph{2018}\label{section-43}

\textbf{\emph{Pawan
Bharadwaj}},
Laurent
Demanet,
Aimé
Fournier,
Systems
and
methods
for
focused
blind
deconvolution,
U.S.
Patent
Application
No.
16/407,884.

\paragraph{2017}\label{section-44}

\textbf{\emph{Pawan
Bharadwaj}},
Laurent
Demanet,
Aimé
Fournier,
Deblending
random
seismic
sources
via
independent
component
analysis,
U.S.
Provisional
Patent
Application
62481767.

\cvsect{arXiv.org}\label{arxiv.org}

\paragraph{2017}\label{section-45}

\textbf{\emph{Pawan
Bharadwaj}},
Wim
Mulder,
and Guy
Drijkoningen,
A
parameterization
analysis
for
acoustic
full-waveform
inversion
of
sub-wavelength
anomalies.
\href{https://github.com/pawbz/pawbz.github.io/blob/src/pdf/papers/Param_arXiv.pdf}{{[}pdf{]}}

\begin{minipage}[t]{0.2\textwidth}
	

	\cvsect{Languages}
	\vspace{\baselineskip} % Required for vertically aligning minipages
	
	\textbf{English} --- proficient\\
	\textbf{Telugu} --- native\\
	\textbf{Hindi} --- proficient
\end{minipage}
\hfill
\begin{minipage}[t]{0.8\textwidth}
	
	\cvsect{References}
%	\vspace{\baselineskip} % Required for vertically aligning minipages
	
	\begin{itemize}
			\setlength\itemsep{0.5pt}
		\item Wim Mulder,  Professor of Geophysical Imaging Methods, Delft University of Technology.
		\item Gerard Schuster,  Professor of Geophysics, King Abdullah University of Science and Tech.
		\item Laurent Demanet, Professor of Mathematics, Massachusetts Institute of Technology.
 
	\end{itemize}
\end{minipage}

%----------------------------------------------------------------------------------------

\end{document}
